\documentclass[11pt]{article}
\usepackage[T1]{fontenc}
\usepackage{lmodern,microtype}
\usepackage[margin=1in]{geometry}
\usepackage{amsmath,amssymb,amsthm,mathtools,booktabs,array,xcolor}
\usepackage[colorlinks=true,linkcolor=blue!45!black,citecolor=blue!45!black,urlcolor=blue!45!black]{hyperref}
\usepackage[nameinlink,noabbrev]{cleveref}
\setlength{\emergencystretch}{2em}
\newtheorem{theorem}{Theorem}[section]
\newtheorem{proposition}[theorem]{Proposition}
\newtheorem{corollary}[theorem]{Corollary}
\newtheorem{lemma}[theorem]{Lemma}
\theoremstyle{remark}\newtheorem{remark}[theorem]{Remark}
\newcommand{\E}{\mathbb E}\newcommand{\PP}{\mathbb P}
\newcommand{\Var}{\operatorname{Var}}\newcommand{\Cov}{\operatorname{Cov}}
\newcommand{\dd}{\,\mathrm d}\newcommand{\ind}{\mathbf1}
\newcommand{\fall}[2]{(#1)_{#2}}\newcommand{\rise}[2]{(#1)^{\overline{#2}}}
\title{Capture--Recapture in a Sampling-Kernel Framework\\
\large Exact recovery, estimator corrections, transfer, and calibration}
\author{Christopher D. Long}
\date{November 17, 2026\\Research draft, version 0.1}
\begin{document}\maketitle
\begin{abstract}
We develop the capture--recapture component of a partial-population inference
framework. Fixed-size overlap sampling gives reciprocal-factorial recovery,
an exact bias identity for Chapman, and an explicit classification of signed
transfers when population size is unknown. Tail inversion yields exact
confidence intervals with sharp one-sided monotonicity properties. Sampling
with replacement yields Bailey's different correction. Repeated captures give weighted estimating equations and exact test
martingales. Inverse sampling yields complete sufficient order statistics and
connects directly to the German tank experiment. Independent capture histories
recover the counting-measure cumulant geometry, Horvitz--Thompson weighting,
and finite-occasion Chao-type moment bounds. Proofs and rational verification
separate estimator identities, calibration, and optimality in specified
classes. Classical formulas are attributed at use; detailed comparisons with
prior work appear near the end.
\end{abstract}

\section{Experiments and targets}
The organizing objects are a sampling law, a population target, and an
observation-preserving transfer, as in the companion papers
\cite{LongDiscovery,LongTransfer,LongCumulants}. Here the population is closed:
there are no births, deaths, migration, lost marks, or erroneous matches.
Unless explicitly stated otherwise, a capture is a uniform subset and marking
does not change subsequent selection. Labels record identity only. Their
numeric magnitudes carry no information about population size.

First fix positive integers $a,b$, an unknown integer $N\ge\max(a,b)$,
and a marked set of size $a$. An independent uniform subset of size $b$
produces overlap $J$. Conditional on the marked set,
\begin{equation}\label{eq:kernel}
 P_N^{a,b}(j)=\frac{\binom aj\binom{N-a}{b-j}}{\binom Nb},
 \qquad 0\le j\le m:=\min(a,b).
\end{equation}
Binomial coefficients are zero outside their usual nonnegative range.
The observed union imposes $N\ge a+b-j$. Random uniform marking followed by
independent uniform recapture gives the same law. Only the second-stage
uniformity is needed when the marked set is conditioned upon.

This parameterization differs from a hypergeometric experiment with known
population size and unknown marked count. In particular, the polynomial
recovery theorem for that experiment cannot be applied by renaming its
unknown parameter. We write $\fall{x}{r}=x(x-1)\cdots(x-r+1)$ and
$\rise{x}{r}=x(x+1)\cdots(x+r-1)$, with both empty products equal to one.

\section{Exact recovery and the two-sample estimators}
\begin{theorem}[Reciprocal-factorial recovery]\label{thm:range}
For $0\le r\le m$,
\begin{equation}\label{eq:factorial}
 \E_N\fall{J}{r}=\frac{\fall{a}{r}\fall{b}{r}}{\fall{N}{r}}.
\end{equation}
The expectations of all finite-valued statistics of $J$ form exactly
\[
 \mathcal R_{a,b}=\operatorname{span}\{1,\fall{N}{1}^{-1},\ldots,
 \fall{N}{m}^{-1}\}
\]
as functions on the whole parameter space. Each target in this space has a
unique unbiased statistic of $J$. In particular, $N$ has no finite-valued
unbiased estimator based on $J$ over this unbounded parameter space.
\end{theorem}
\begin{proof}
Expand $\fall{J}{r}$ as a sum over ordered distinct marked $r$-tuples.
Each tuple is in the recapture subset with probability
$\fall{b}{r}/\fall{N}{r}$, giving \eqref{eq:factorial}.
The $m+1$ functions $\fall{j}{r}$ form a triangular basis for statistics on
$\{0,\ldots,m\}$. The resulting reciprocal-factorial functions are linearly
independent: multiply any relation by $\fall{N}{m}$. The terms become
polynomials of distinct degrees $m,m-1,\ldots,0$, and a polynomial vanishing
at every sufficiently large integer is zero. This proves the range and
uniqueness statements. Finally, every finite-valued statistic on this finite
observation space has bounded expectation; $N$ is unbounded.
\end{proof}
The uniqueness is within the overlap experiment. Independent randomization
cannot improve squared-error risk: conditional averaging given $J$ removes
its extra variance whenever that risk is finite.

The Lincoln--Petersen ratio \cite{Chapman1951,Dettloff2023} is
\begin{equation}\label{eq:LP}
 \widehat N_{\rm LP}=ab/J,\qquad J>0.
\end{equation}
It inverts the unbiased estimator $J/(ab)$ of $1/N$; reciprocation does not
preserve unbiasedness. If the zero-overlap value is defined as $+\infty$,
its expectation and squared-error risk are infinite whenever $N\ge a+b$.
Any finite replacement at zero defines a different estimator, whose bias
and risk must include that event.

\begin{proposition}[Integer likelihood maximizers]\label{prop:mle}
For observed $j>0$, the likelihood maximizers are $\lfloor ab/j\rfloor$
and, if $ab/j$ is an integer, also $ab/j-1$, retaining only values at least
$a+b-j$. For $j=0$, the likelihood increases to its supremum as
$N\to\infty$ and has no finite maximizer.
\end{proposition}
\begin{proof}
For consecutive feasible values with $N>a+b-j$,
\[
 \frac{P_N^{a,b}(j)}{P_{N-1}^{a,b}(j)}
 =\frac{(N-a)(N-b)}{N(N-a-b+j)}.
\]
It is at least one exactly when $Nj\le ab$. Moreover
$ab/j\ge a+b-j$ because $(a-j)(b-j)\ge0$.
For $j=0$, the same comparison is strict for consecutive feasible values,
and the probability tends to one.
\end{proof}

\begin{theorem}[Chapman and the reciprocal-shift identity]\label{thm:chapman}
Chapman's estimator \cite{Chapman1951} satisfies
\begin{align}
 \widehat N_{\rm C}&=\frac{(a+1)(b+1)}{J+1}-1,\label{eq:Chapman}\\
 \E_N\widehat N_{\rm C}
 &=N-(N-a-b)P_N^{a,b}(0).\label{eq:Chapmanbias}
\end{align}
More generally, for every integer $k\ge1$,
\begin{align}\label{eq:shift}
 \E_N\frac{\rise{a+1}{k}\rise{b+1}{k}}{\rise{J+1}{k}}
 &=\rise{N+1}{k}-B_k(N),\\
 B_k(N)&=\frac{\rise{b+1}{k}}{\binom Nb}
 \sum_{h=0}^{k-1}\binom{a+k}{h}\binom{N-a}{b+k-h}\ge0.\nonumber
\end{align}
Thus Chapman is unbiased when $a+b\ge N$, and strictly downward biased
when $a+b<N$. Its reported value is at least the observed union size.
\end{theorem}
\begin{proof}
Use
\[
 \frac{\binom aj}{\rise{j+1}{k}}
 =\frac{\binom{a+k}{j+k}}{\rise{a+1}{k}}.
\]
After setting $h=j+k$, Vandermonde's convolution completes the sum to
$\binom{N+k}{b+k}$; the omitted indices are exactly $0\le h<k$.
The completed term equals $\rise{N+1}{k}$. At $k=1$ the remainder is
\[
 B_1(N)=\frac{(b+1)\binom{N-a}{b+1}}{\binom Nb}
       =(N-a-b)P_N^{a,b}(0),
\]
with both sides zero when $N<a+b$. Finally,
$\widehat N_{\rm C}-(a+b-j)=(a-j)(b-j)/(j+1)\ge0$.
\end{proof}
The correction is a kernel boundary term, not a proof of global unbiasedness.
For any finite rule $h(j)$ the exact bias and risk are computable as
\begin{equation}\label{eq:risk}
 \operatorname{Bias}_N(h)=\sum_jP_N^{a,b}(j)h(j)-N,
 \qquad R_N(h)=\sum_jP_N^{a,b}(j)(h(j)-N)^2.
\end{equation}
In particular, \eqref{eq:risk} supplies exact Chapman variance and MSE;
no approximate plug-in variance formula is needed for the finite-design audit.

\section{Changing the second-sample design}
Suppose instead that the second stage consists of $b$ independent draws
\emph{with replacement}. Repeated sightings of the same marked unit count
again. Then $J\sim\operatorname{Bin}(b,a/N)$, with $0\le J\le b$ even
when $b>a$.
\begin{proposition}[Bailey's binomial correction]\label{prop:bailey}
The estimator attributed to Bailey \cite{Bailey1951,Bailey1952} obeys
\[
 \widehat N_{\rm B}=\frac{a(b+1)}{J+1},\qquad
 \E_N\widehat N_{\rm B}
 =N\left[1-\left(1-\frac aN\right)^{b+1}\right].
\]
The unbiasedly estimable functions from $J$ are exactly polynomials in
$1/N$ of degree at most $b$.
\end{proposition}
\begin{proof}
Put $p=a/N$. The identity
$\binom bj/(j+1)=\binom{b+1}{j+1}/(b+1)$ gives
$\E[1/(J+1)]=[1-(1-p)^{b+1}]/[(b+1)p]$.
The factorial moments are $\E\fall{J}{r}=\fall{b}{r}(a/N)^r$;
the same basis argument as in \cref{thm:range} proves the range statement.
\end{proof}
Using Chapman in this binomial design changes the bias calculation. An
approximation of a hypergeometric law by a binomial law is not an exact
transfer assertion. Exact intervals use the binomial tails on $N\ge a$.
Their lower and upper endpoints are obtained from the same strict inequalities
as \cref{thm:limits}: the binomial count decreases stochastically as $a/N$
decreases. The upper endpoint at zero recaptures is infinite. The code
implements this design separately, without imposing the without-replacement
restriction $J\le a$ or $b\le N$.

\section{Transfer between unknown-size overlap experiments}
All experiments in this section have common parameter space
$N\ge\max(a,b,c,d)$, with no upper bound.
\begin{theorem}[Explicit signed-transfer classification]\label{thm:transfer}
A normalized signed matrix $S$ satisfying
$P_N^{a,b}S=P_N^{c,d}$ for all such $N$ exists if and only if
$h:=\min(c,d)\le m:=\min(a,b)$. When it exists it is unique and
\begin{equation}\label{eq:transfer}
 S(j,y)=\sum_{r=y}^{h}(-1)^{r-y}\binom ry\binom jr
 \frac{\fall{c}{r}\fall{d}{r}}{\fall{a}{r}\fall{b}{r}},
 \quad 0\le j\le m,\quad0\le y\le h.
\end{equation}
It is a positive transfer exactly when all these entries are nonnegative.
Its unique row-variation amplification is $\max_j\sum_y|S(j,y)|$.
\end{theorem}
\begin{proof}
By \cref{thm:range}, source and target column spaces are the reciprocal-
factorial spaces of degrees $m$ and $h$, respectively. Their independence
continues to hold on the common infinite tail of integers, so target range
inclusion requires $h\le m$.
For this case define each signed row through its generating polynomial
\[
 \sum_y S(j,y)z^y=
 \sum_{r=0}^h\binom jr
 \frac{\fall{c}{r}\fall{d}{r}}{\fall{a}{r}\fall{b}{r}}(z-1)^r.
\]
It has value one at $z=1$. Taking source expectations gives the target
factorial-moment generating polynomial, hence the target law.
Coefficient extraction gives \eqref{eq:transfer}. Full source column rank
implies uniqueness. Positivity and the variation formula follow directly.
\end{proof}
For $c\le a,d\le b$, there is a direct positive construction: retain a
uniform $c$-subset of the marked set, then a uniform $d$-subset of the
recapture set. Conditional on $J=j$, first draw
$H\sim\operatorname{Hypergeom}(a,j,c)$ and then
$Y\mid H\sim\operatorname{Hypergeom}(b,H,d)$. Uniqueness identifies this
composition with \eqref{eq:transfer}.

Existence alone does not preserve calibration. For $(a,b)=(2,2)$ and
$(c,d)=(1,3)$, the unique matrix is
\[
 S=\begin{pmatrix}1&0\\1/4&3/4\\-1/2&3/2\end{pmatrix}.
\]
Its amplification is two and it is not positive. The hypotheses of the
positive/dominated calibration transfer results in \cite{LongTransfer}
cannot be replaced by signed range inclusion. On a finite restricted grid
of $N$ values the range criterion may differ; the infinite parameter space
in this theorem is essential.

\section{Exact fixed-size confidence limits}
Write $F_N(j)=P_N(J\le j)$ and $G_N(j)=P_N(J\ge j)$, extending the
probabilities to all $N\ge\max(a,b)$ with impossible cells assigned zero.
For $0<\alpha<1$ define
\begin{equation}\label{eq:limits}
 L_\alpha(j)=\min\{N:F_N(j)>\alpha\},\qquad
 U_\alpha(j)=\sup\{N:G_N(j)>\alpha\}.
\end{equation}
Both sets incorporate the actual support: retained $N$ satisfy
$N\ge a+b-j$. The upper endpoint is an integer or $+\infty$.
\begin{theorem}[Coverage and one-sided sharpness]\label{thm:limits}
The limits in \eqref{eq:limits} are nonincreasing in $j$, and
\[
 P_N\{L_\alpha(J)\le N\}\ge1-\alpha,
 \qquad P_N\{N\le U_\alpha(J)\}\ge1-\alpha.
\]
Among nonincreasing, nonrandomized integer lower limits with this coverage,
$L_\alpha$ is pointwise largest. Among nonincreasing upper limits, allowing
$+\infty$, $U_\alpha$ is pointwise smallest.
For positive $\alpha_L,\alpha_U$ with sum $\alpha$, the set
\[
 C(j)=\{N:F_N(j)>\alpha_L,\ G_N(j)>\alpha_U\}
 =[L_{\alpha_L}(j),U_{\alpha_U}(j)]\cap\mathbb Z
\]
has coverage at least $1-\alpha$. No shortest two-sided interval assertion
is made.
\end{theorem}
\begin{proof}
Add unmarked individuals to the population and assign independent continuous
priorities to all units. Select the $b$ smallest priorities. Adding unmarked
units cannot increase the number marked in that subset. Hence $F_N(j)$ is
nondecreasing in $N$ and $G_N(j)$ is nonincreasing. At fixed $N$, $F_N(j)$
increases and $G_N(j)$ decreases with $j$, proving the endpoint monotonicity.
Also $F_N(j)\to1$ as $N\to\infty$, so $L_\alpha$ exists.
The elementary discrete tail inequalities
$P_N\{F_N(J)\le\alpha\}\le\alpha$ and
$P_N\{G_N(J)\le\alpha\}\le\alpha$ give coverage and the union bound.

If a competing nonincreasing lower limit $V$ has
$V(j)>L_\alpha(j)$, take $N=L_\alpha(j)$. On $J\le j$ it overestimates,
with probability $F_N(j)>\alpha$. For a smaller upper limit at a finite
$U_\alpha(j)$, take $N=U_\alpha(j)$ and use the event $J\ge j$.
At $j=0$, $P_N(J=0)\to1$, so any finite upper value fails coverage for
large $N$. These arguments prove the pointwise assertions.
\end{proof}
In particular $U_\alpha(0)=+\infty$, whereas $U_\alpha(j)<\infty$ for
$j\ge1$, since $G_N(j)\le ab/(Nj)$. At $N=a+b-j$, $G_N(j)=1$, so the
upper acceptance set is nonempty. Strict inequalities are needed for the
stated pointwise minimality/maximality. The implementation brackets the
monotone crossings by doubling and uses rational binary search; it never
replaces infinity by an arbitrary population ceiling.

For a different exact calibration, choose any predetermined probability
mass function $q(j)$ on the observation space. For each candidate $N$,
$e_N(j)=q(j)/P_N(j)$ on its positive cells has expectation at most one.
Impossible observed cells reject the candidate. Markov's inequality gives
$\{N:e_N(J)<1/\alpha\}$ coverage at least $1-\alpha$.
This construction is valid independently of any selected point estimator;
it need not produce an interval or a shortest confidence set.

\section{Repeated captures and sequential calibration}
Before occasion $t$ let $a_t$ be the number marked, choose a predictable
sample size $b_t$, and obtain a uniform subset of that size. Return all
sampled individuals and mark every new one. Then
\begin{equation}\label{eq:conditional}
 J_t\mid\mathcal F_{t-1}\sim P_N^{a_t,b_t},\qquad
 a_{t+1}=a_t+b_t-J_t,\qquad
 \E_N[J_t\mid\mathcal F_{t-1}]=a_tb_t/N.
\end{equation}
Zero marked counts are permitted at an initial occasion. Chosen actions
must be feasible under the true population. The policy is a known function
of the past, not a function of the unknown $N$.

\begin{proposition}[Weighted estimating functions]\label{prop:weighted}
For predictable finite weights $w_t$, the increments
$w_t(J_t-a_tb_t/N)$ have conditional mean zero whenever integrable.
At a fixed horizon the estimating equation has ratio solution
\begin{equation}\label{eq:weighted}
 \widehat N_w=\frac{\sum_t w_ta_tb_t}{\sum_t w_tJ_t}
\end{equation}
when its denominator is nonzero. Weights $w_t=1$ yield Schnabel
\cite{Schnabel1938}; weights $w_t=a_t$ yield Schumacher--Eschmeyer
\cite{Schumacher1943,Dettloff2023}.
\end{proposition}
\begin{proof}
Multiply the last equality of \eqref{eq:conditional} by the predictable
weight, then sum. Solving the resulting equation gives \eqref{eq:weighted}.
\end{proof}
A zero-mean estimating function does not imply an unbiased ratio solution.
Nor does it give an optional-stopping theorem without additional conditions.
For a fixed size sequence, an exact risk/distribution calculation propagates
states $(a,A,B)$, where $A=\sum w_ta_tb_t$ and $B=\sum w_tJ_t$:
\[
 (a,A,B)\longmapsto(a+b-j,A+wab,B+wj)
 \quad\hbox{with probability }P_N^{a,b}(j).
\]
The terminal distribution retains any mass at an infinite ratio. It computes
finite losses, conditional risks, or the risk of an explicitly specified
finite fallback; it cannot silently discard zero-recapture trajectories.

\begin{theorem}[Exact test martingales]\label{thm:martingale}
For any finite predictable real directions $\lambda_t$, put
\[
 \psi_{N,t}(\lambda)=\log\sum_jP_N^{a_t,b_t}(j)e^{\lambda j},\qquad
 E_{N,T}=\exp\sum_{t\le T}\{\lambda_tJ_t-\psi_{N,t}(\lambda_t)\}.
\]
Under the candidate's true law this is a nonnegative martingale starting at
one. Consequently
\[
 P_N\{\exists T:E_{N,T}\ge1/\alpha\}\le\alpha.
\]
The running sets $\{N:\max_{s\le T}E_{N,s}<1/\alpha\}$ form a confidence
sequence. Observed prefixes impossible under a candidate may reject it.
Predetermined mixtures of such martingales are valid as well.
\end{theorem}
\begin{proof}
Conditional on the past, the next multiplicative factor has expectation one.
Every finite-step normalizer is positive and finite because its observation
support is finite. Iteration gives integrability and expectation one.
For the first crossing time $\tau$, the bounded-stopping identity gives
$1=\E E_{N,\tau\wedge T}\ge\alpha^{-1}P(\tau\le T)$.
Let $T\to\infty$. Mixtures follow by nonnegativity and Tonelli's theorem.
\end{proof}
This is the standard test-martingale principle \cite{Howard2021}; martingale
estimating functions in recapture studies are classical \cite{Lloyd1994}.
The confidence sequence permits arbitrary subsequent stopping. Selecting
$\lambda_t$ using $J_t$ or later outcomes does not satisfy the theorem.

\section{Inverse sampling and the tank equivalence}
Fix $1\le r\le a$ before sampling. Draw without replacement from a random
permutation of the population until the $r$th marked unit appears; call the
number of draws $T_r$. The complete observed binary history ends at that
marked draw. Its law is
\begin{equation}\label{eq:inverse}
 P_N(T_r=t)=\frac{\binom{t-1}{r-1}\binom{N-t}{a-r}}{\binom Na},
 \qquad r\le t\le N-a+r.
\end{equation}
This is the inverse hypergeometric design studied in the classical recapture
literature \cite{Bailey1951,Chapman1951} and by
P\'erez-Salvador \cite{Perez2026}.

\begin{theorem}[Order-statistic recovery]\label{thm:inverse}
In this experiment $T_r$ is complete and sufficient for the binary history,
and
\begin{align}
 \E_N T_r&=\frac{r(N+1)}{a+1},\\
 \widehat N_r&=\frac{a+1}{r}T_r-1,\qquad
 \Var_N(\widehat N_r)=\frac{(N+1)(N-a)(a+1-r)}{r(a+2)}.\label{eq:invvariance}
\end{align}
Thus $\widehat N_r$ is minimum-variance unbiased for $N$ among estimators
of this binary history. At $r=a$ the experiment based on $T_a$ is exactly
the maximum-statistic experiment for a uniform $a$-subset of
$\{1,\ldots,N\}$, giving the German tank formula.
\end{theorem}
\begin{proof}
The $a$ marked positions are a uniform $a$-subset, proving \eqref{eq:inverse}.
Its $k=a+1$ empty gaps form a uniform weak composition of $L=N-a$.
For distinct gap indices, coefficient extraction from $(1-z)^{-k}$ gives
\[
 \E G_i=L/k,\quad
 \E\fall{G_i}{2}=\frac{2L(L-1)}{k(k+1)},\quad
 \E G_iG_j=\frac{L(L-1)}{k(k+1)}.
\]
For example the second numerator comes from
$2z^2(1-z)^{-3}(1-z)^{-(k-1)}$ divided by
$\binom{L+k-1}{k-1}$. Since
$T_r=r+G_0+\cdots+G_{r-1}$, these moments give the stated mean and variance.

Conditional on $T_r=t$, the $r-1$ earlier marked positions are a uniform
subset of $\{1,\ldots,t-1\}$, independent of $N$, proving sufficiency.
If $\E_N h(T_r)=0$ for every $N\ge a$, start at $N=a$ to obtain $h(r)=0$.
At $N=a+s$, support is $r,\ldots,r+s$ with positive mass at $r+s$;
induction gives $h(r+s)=0$. This proves completeness. Conditional averaging
and completeness establish minimum variance. The case $r=a$ is the maximum
of the uniform subset, with no approximation or change in its law.
\end{proof}
Let $H_N(t)=P_N(T_r\le t)$ and $K_N(t)=P_N(T_r\ge t)$.
An exact equal-tail interval at a preassigned $r$ is
\begin{equation}\label{eq:inverseCI}
 \{N\ge a+t-r:H_N(t)>\alpha/2,\ K_N(t)>\alpha/2\}.
\end{equation}
Its coverage follows from the same discrete-tail argument as
\cref{thm:limits}. It is an integer interval: adding unmarked units to a
random permutation cannot move a marked position earlier, so $H_N(t)$
decreases and $K_N(t)$ increases in $N$.
For computation, $H_N(t)=P\{\operatorname{Hypergeom}(N,a,t)\ge r\}$ when
$t\le N$, with $H_N(t)=1$ when $t\ge N$.
Both endpoints are finite: $K_N(t)\to1$ and
$H_N(t)\le at/(Nr)\to0$ as $N\to\infty$.
No fixed-$r$ coverage claim is made for a data-selected $r$.

For arbitrary stopping while drawing without replacement, the marked
probability at draw $s$ conditional on the preceding $k$ marked draws is
$p_{N,s}=(a-k)/(N-s+1)$. For predictable $z_s>0$ the factors
\[
 \frac{z_s^{X_s}}{1-p_{N,s}+p_{N,s}z_s},\qquad X_s\in\{0,1\},
\]
have conditional expectation one. Their products provide a confidence
sequence by \cref{thm:martingale}'s argument, with impossible prefixes
excluded and the process frozen when the candidate population is exhausted.
This supplies the stopping-valid alternative to reusing \eqref{eq:inverseCI}.

With replacement, stopping at $r\ge1$ marked \emph{draws} instead yields
\[
 P_N(T=t)=\binom{t-1}{r-1}(a/N)^r(1-a/N)^{t-r},\quad t\ge r.
\]
The sum of $r$ independent geometric waiting times gives
\[
 \E_N T=rN/a,\qquad \widehat N_{\rm IB}=aT/r,
 \qquad\Var_N(\widehat N_{\rm IB})=N(N-a)/r.
\]
There is no restriction $r\le a$ because repeats count. Exact intervals use
the negative-binomial tails, or
$P_N(T\le t)=P\{\operatorname{Bin}(t,a/N)\ge r\}$.
This is a different experiment from the without-replacement order statistic;
we make no completeness claim for its discrete parameter subfamily.

\section{Independent histories, inclusion weighting, and heterogeneity}
For $K$ occasions let $W_i\in\{0,1\}^K$ be the history of individual $i$.
Conditional on fixed catchability vectors $q_i\in[0,1]^K$, assume histories
are independent across individuals and captures are independent across
occasions within an individual. Put
\[
 \nu=\sum_{i=1}^N\delta_{q_i},\quad
 \pi_h(q)=\prod_{s=1}^K q_s^{h_s}(1-q_s)^{1-h_s},\quad
 C_h=\sum_i\ind\{W_i=h\}.
\]
Here $N=\nu([0,1]^K)$, $S=\sum_{h\ne0}C_h$, and the realized unseen count
is $U=C_0=N-S$. The population measure has unknown total mass; it is not
the normalized abundance measure used for species probabilities.

\begin{theorem}[Counting-measure cumulant calibration]\label{thm:history}
For every finite real vector $(\lambda_h)_{h\ne0}$ and real $\eta$,
\begin{equation}\label{eq:cumulant}
 \log\E\exp\left(\sum_{h\ne0}\lambda_h C_h+\eta U\right)
 =\int\log\left(\pi_0(q)e^\eta+
       \sum_{h\ne0}\pi_h(q)e^{\lambda_h}\right)\dd\nu(q).
\end{equation}
Fix finitely many directions and nonnegative mixture weights summing to one.
For candidate pairs $(\nu,u)$, the corresponding mixture of normalized
exponentials has expectation one at the true population and realized target.
Its strict sublevel set at $1/\alpha$ is jointly convex over finite positive
measures and $u\in\mathbb R$, and contains $(\nu,U)$ with probability at
least $1-\alpha$. Intersecting with $\nu([0,1]^K)=S+u$ preserves convexity
and coverage. Projection gives intervals (possibly unbounded or empty);
restricting to integer-realizable populations preserves coverage but can
destroy convexity.
\end{theorem}
\begin{proof}
Independence factors the expectation into $N$ positive finite one-individual
transforms. Taking logarithms gives \eqref{eq:cumulant}. For each direction,
subtracting the right side normalizes its exponential to expectation one.
A finite nonnegative mixture retains this property, and Markov's inequality
gives coverage. At fixed observed counts its exponent is affine in $(\nu,u)$,
so the mixture is convex. The additional equality is affine and holds at
$(\nu,U)$ for every outcome. A convex subset projects to an interval in a
real coordinate. Actual finite populations satisfy the equality and are a
subset of the relaxed class, giving the last assertion.
\end{proof}
The integrand is bounded on the compact catchability cube for finite
directions, so every finite measure is admissible. This is the independent-
summand identity of \cite{LongCumulants} in counting-measure coordinates;
it does not establish new independence in a fixed-size capture design.
If the $q_i$ themselves are random and integrated out, the marginal transform
must also be integrated out. In particular, for iid $q_i\sim F$, it is
$N\log\int(\pi_0e^\eta+\sum_{h\ne0}\pi_he^{\lambda_h})\dd F$, not the
integral of the logarithm against $NF$.

For two homogeneous Bernoulli occasions with probabilities in $(0,1)$,
conditioning on the two sample totals $(a,b)$ makes the two subsets independent
and uniform. To see this, each size-$a$ first subset has the same probability
$q_1^a(1-q_1)^{N-a}$, and likewise for the second; the events factor across
occasions. Their conditional overlap law is \eqref{eq:kernel}, with the
nuisance probabilities canceled. Unequal catchabilities give weighted subsets
instead, and that cancellation cannot be assumed.

\begin{proposition}[Known inclusion probabilities]\label{prop:HT}
In any sampling design let $I_i$ indicate eventual inclusion and suppose
$\rho_i=P(I_i=1)>0$ is known for every individual. Then the
Horvitz--Thompson statistic \cite{HT1952}
\[
 \widehat N_{\rm HT}=\sum_{i:I_i=1}\rho_i^{-1}
\]
is unbiased. If $\rho_{ij}=P(I_i=I_j=1)$, its variance is
\[
 \sum_i\frac{1-\rho_i}{\rho_i}
 +2\sum_{i<j}\frac{\rho_{ij}-\rho_i\rho_j}{\rho_i\rho_j}.
\]
For independent histories $\rho_i=1-\pi_0(q_i)$ and the second sum vanishes.
\end{proposition}
\begin{proof}
Use $\E I_i=\rho_i$, $\Var(I_i)=\rho_i(1-\rho_i)$, and expand the variance
of the weighted sum. Independence is needed only for the final simplification.
\end{proof}
In classical capture--recapture the inclusion probabilities are unknown.
Replacing them by estimates need not preserve unbiasedness or coverage.

\begin{proposition}[Finite-occasion Chao-type moment bound]\label{prop:chao}
Suppose individual $i$ has the same detection probability $q_i$ on each of
$K\ge2$ independent occasions. Let $f_j$ count individuals seen on exactly
$j$ occasions and $m_j=\E f_j$. Then
\begin{equation}\label{eq:chao}
 m_0\ge\frac{K-1}{2K}\frac{m_1^2}{m_2}\qquad(m_2>0).
\end{equation}
For independent Poisson detection counts, the analogous bound is
$m_0\ge m_1^2/(2m_2)$. Both are inequalities between population moments,
not finite-sample confidence assertions for ratios of observed frequencies.
\end{proposition}
\begin{proof}
Write $A_i=(1-q_i)^{K/2}$ and
$B_i=q_i(1-q_i)^{(K-2)/2}$, with endpoint values interpreted continuously.
Then $\sum A_i^2=m_0$, $K\sum A_iB_i=m_1$, and
$\binom K2\sum B_i^2=m_2$. Cauchy--Schwarz proves \eqref{eq:chao}.
For Poisson means $x_i$, use $A_i=e^{-x_i/2}$ and $B_i=x_ie^{-x_i/2}$.
\end{proof}
The associated incidence plug-in rule, when $f_2>0$, is
\[
 \widehat N_{\mathrm{Chao},K}
 =S+\frac{K-1}{2K}\frac{f_1^2}{f_2}.
\]
For Poisson detection the corresponding rule is $S+f_1^2/(2f_2)$.
These ratios need not be unbiased, and neither is a samplewise lower bound
or a confidence bound. Our incidence implementation returns $+\infty$ when
$f_2=0<f_1$; when $f_1=f_2=0$ it returns $S$ as an explicit point-estimate
convention only. Other regularizations define different procedures requiring
their own risk analysis. Exact laws can be computed by enumerating the
independent histories, retaining zero-doubleton outcomes.

The binomial coefficient in \eqref{eq:chao} is retained exactly; a Poisson
approximation would remove it. Chao's capture--recapture work and its moment
connection are prior art \cite{Chao1987}. Moment duality and additional
constraints can be applied to these kernels exactly as in
\cite{LongDiscovery}, but the exponential discovery envelopes are not
asserted to remain the extremizers for a different kernel family.

\begin{proposition}[Undetectable-population obstruction]\label{prop:obstruction}
If zero catchability vectors are allowed, adding any number of such
individuals leaves the law of all observed nonzero histories unchanged.
If all coordinates are positive but arbitrarily small, no uniformly valid
upper population bound at error $\alpha<1$ can be almost surely finite
under every finite baseline population.
\end{proposition}
\begin{proof}
The first assertion follows from $W_i=0$ surely for an added zero vector.
For the second, fix any baseline population and suppose its reported upper
bound is finite almost surely. Choose $\delta>0$ with $1-2\delta>\alpha$.
There is a finite $B$ such that the baseline reports a value at most $B$
with probability exceeding $1-\delta$. Add $M$ individuals so total size
exceeds $B$, assigning them positive catchabilities small enough that their
probability of any detection is less than $\delta$ (the union bound over
$MK$ Bernoulli captures suffices). Couple the baseline histories identically.
With probability greater than $1-2\delta$ the enlarged population produces
the same observations and an upper value at most $B$, violating coverage.
\end{proof}
A known lower detection bound or a specified heterogeneity model can remove
this particular obstruction. Neither is implied by a large observed sample.

\section{Relation to prior work}\label{sec:prior}
The estimators named here are classical, and their presence in one framework
is not a claim of first unification. Chapman \cite{Chapman1951} and Bailey
\cite{Bailey1951,Bailey1952} treat bias corrections and distinct recapture
designs. Schnabel \cite{Schnabel1938} and Schumacher--Eschmeyer
\cite{Schumacher1943} supply repeated-capture ratios. Horvitz--Thompson
\cite{HT1952} supplies general inclusion weighting. Dettloff
\cite{Dettloff2023} compares these simple estimators and proposes further
bias adjustment; its empirical comparisons are not invoked as universal
risk or coverage theorems here.

The connection between heterogeneous capture--recapture, unseen species,
and moment recovery was already explicit in Chao \cite{Chao1987}, including
a comparison with Harris's moment problem. Lloyd \cite{Lloyd1994} studies
martingale estimation, and Chao, Yip, Lee, and Chu \cite{Chao2001} explicitly
propose a unified estimating-function and sample-coverage approach. These
are substantive antecedents to the present framework. The finite-occasion
coefficient in \cref{prop:chao} follows directly from the exact binomial
kernel; the Poisson version must not silently replace that kernel.

Inverse sampling is also classical. P\'erez-Salvador \cite{Perez2026}
explicitly treats the without-replacement estimator and variance in
\eqref{eq:invvariance}, together with approximate intervals and an adaptive
stability rule. Here the preassigned-$r$ interval uses the exact law, and
arbitrary stopping is justified separately by a test martingale. The
fixed-$r$ proof alone supplies no validity at a selected stopping index.
The tank equivalence is an elementary order-statistic identification.

The contribution of this module is an explicit interface: reciprocal-factorial
range characterization, a fully specified signed-transfer matrix, boundary
terms for reciprocal shifts, one-sided monotone limits with exact witnesses,
and the connection of repeated and independent-history experiments to the
companion calibration machinery. Each statement is proved here; a first-
appearance claim for these algebraic consequences would require additional
literature comparison. In particular, the signed-transfer classification is
an application of the established range principle, not a replacement for
Blackwell comparison theory. There is no universal shortest-interval,
minimax, or heterogeneity-robustness claim.

\section{Verification and incorporation}
The module \texttt{src/partial\_population/capture\_recapture.py} uses exact
rational arithmetic for finite probabilities, corrections, transfers,
confidence endpoints, and test-martingale paths. Its verifier compares
combinatorial kernels with explicit subsets and permutations, checks the
factorial and reciprocal-shift identities, enumerates confidence coverage,
and checks martingale expectations and crossing probabilities under
predictable choices. An independent history enumeration checks the product
transform. Numerical population caps are not used to truncate confidence
endpoints. Exact arithmetic avoids roundoff but is not optimized for large
sample sizes.

The discovery companion uses this module as its capture--recapture
specialization; the transfer companion distinguishes its unknown-$N$ range
from the known-population polynomial range; the cumulant companion uses
its counting-measure formulation for independent histories. Existing
manuscript theorem numbering is preserved by these cross-references.

This development covers the stated closed-population experiments. Open
populations, tag loss, match error, spatial detection, behavioral changes,
and a sharp bounded-heterogeneity inverse-coupon procedure require additional
models and proofs. This paper neither recovers nor attributes any unavailable
older private tag--retag manuscript.

\begingroup\small
\begin{thebibliography}{99}\raggedright
\setlength{\itemsep}{3pt}\setlength{\parskip}{0pt}
\bibitem{Bailey1951} N. T. J. Bailey.
On estimating the size of mobile populations from recapture data.
\emph{Biometrika} \textbf{38}(3--4) (1951), 293--306.
\href{https://doi.org/10.1093/biomet/38.3-4.293}{doi:10.1093/biomet/38.3-4.293}.
\bibitem{Bailey1952} N. T. J. Bailey.
Improvements in the interpretation of recapture data.
\emph{Journal of Animal Ecology} \textbf{21}(1) (1952), 120--127.
\href{https://doi.org/10.2307/1913}{doi:10.2307/1913}.
\bibitem{Chapman1951} D. G. Chapman.
\emph{Some Properties of the Hypergeometric Distribution with Applications
to Zoological Sample Censuses}. University of California Publications in
Statistics, vol. 1, no. 7. University of California Press, 1951.
\href{https://catalog.hathitrust.org/Record/005761134}{Catalog record}.
\bibitem{Chao1987} A. Chao.
Estimating the population size for capture--recapture data with unequal catchability.
\emph{Biometrics} \textbf{43}(4) (1987), 783--791.
\href{https://doi.org/10.2307/2531532}{doi:10.2307/2531532}.
\bibitem{Chao2001} A. Chao, P. S. F. Yip, S.-M. Lee, and W. Chu.
Population size estimation based on estimating functions for closed
capture--recapture models. \emph{Journal of Statistical Planning and Inference}
\textbf{92}(1--2) (2001), 213--232.
\href{https://doi.org/10.1016/S0378-3758(00)00151-8}{doi:10.1016/S0378-3758(00)00151-8}.
\bibitem{Dettloff2023} K. Dettloff.
Assessment of bias and precision among simple closed population mark--recapture
estimators. \emph{Fisheries Research} \textbf{265} (2023), 106756.
\href{https://doi.org/10.1016/j.fishres.2023.106756}{doi:10.1016/j.fishres.2023.106756}.
\bibitem{Howard2021} S. R. Howard, A. Ramdas, J. McAuliffe, and J. Sekhon.
Time-uniform, nonparametric, nonasymptotic confidence sequences.
\emph{Annals of Statistics} \textbf{49}(2) (2021), 1055--1080.
\href{https://doi.org/10.1214/20-AOS1991}{doi:10.1214/20-AOS1991}.
\bibitem{HT1952} D. G. Horvitz and D. J. Thompson.
A generalization of sampling without replacement from a finite universe.
\emph{Journal of the American Statistical Association} \textbf{47}(260)
(1952), 663--685.
\href{https://doi.org/10.1080/01621459.1952.10483446}{doi:10.1080/01621459.1952.10483446}.
\bibitem{Lloyd1994} C. J. Lloyd.
Efficiency of martingale methods in recapture studies.
\emph{Biometrika} \textbf{81}(2) (1994), 305--315.
\href{https://doi.org/10.1093/biomet/81.2.305}{doi:10.1093/biomet/81.2.305}.
\bibitem{LongDiscovery} C. D. Long.
\emph{Shape-Constrained Discovery-Rate Inversion}. Research draft, version 0.3,
November 17, 2026. Repository: \texttt{manuscripts/discovery\_rate\_inversion.tex}.
\bibitem{LongTransfer} C. D. Long.
\emph{Certified Transfer and Generation in Partial-Population Inference}.
Research draft, version 0.2, November 17, 2026.
Repository: \texttt{manuscripts/transfer\_generation.tex}.
\bibitem{LongCumulants} C. D. Long.
\emph{Additive Cumulants and Computable Fixed-Sample Calibration}.
Research note, version 0.2, November 17, 2026.
Repository: \texttt{research/master\_identity\_and\_deletion.tex}.
\bibitem{Perez2026} B. R. P\'erez-Salvador.
Sequential capture--recapture method.
\emph{Sequential Analysis} \textbf{45}(3) (2026), 417--437.
\href{https://doi.org/10.1080/07474946.2026.2620792}{doi:10.1080/07474946.2026.2620792}.
\bibitem{Schnabel1938} Z. E. Schnabel.
The estimation of the total fish population of a lake.
\emph{American Mathematical Monthly} \textbf{45}(6) (1938), 348--352.
\href{https://doi.org/10.1080/00029890.1938.11990818}{doi:10.1080/00029890.1938.11990818}.
\bibitem{Schumacher1943} F. X. Schumacher and R. W. Eschmeyer.
The estimation of fish populations in lakes and ponds.
\emph{Journal of the Tennessee Academy of Science} \textbf{18} (1943), 228--249.
\end{thebibliography}
\endgroup
\end{document}
